% Source: https://artofproblemsolving.com/community/c3h2038154_problems
% Stable identity: aops-topic-2038154-q01
Reflecting in the \(x\)-axis negates \(y\): \(Q=(3,-2)\). Reflecting in \(y=-x\) sends \((x,y)\) to \((-y,-x)\): \(R=(2,-3)\). Reflecting through the origin negates both: \(S=(-2,3)\).

\begin{asy}
usepackage("amsmath");
size(4.5cm);
defaultpen(fontsize(10pt));
draw((-3.5,0)--(4,0),Arrow(4)); draw((0,-3.8)--(0,3.8),Arrow(4));
draw((-3.5,3.5)--(3.5,-3.5),dashed+gray); label("$y=-x$",(3.5,-3.5),E,gray);
pair P=(3,2), Q=(3,-2), R=(2,-3), Sp=(-2,3);
fill(P--Q--R--Sp--cycle,paleblue); draw(P--Q--R--Sp--cycle,linewidth(0.9));
dot(P);dot(Q);dot(R);dot(Sp);
label("$P(3,2)$",P,E); label("$Q(3,-2)$",Q,E); label("$R(2,-3)$",R,S); label("$S(-2,3)$",Sp,N);
label("\mbox{}",(min(currentpicture,user=true)+max(currentpicture,user=true))/2,black); // keeps the text after the figure black when LaTeX typesets the labels inline
\end{asy}

By the shoelace formula with \(P,Q,R,S\) in order,
\[
\begin{aligned}
2[PQRS]&=\left|(3\cdot(-2)-3\cdot2)+(3\cdot(-3)-2\cdot(-2))+(2\cdot3-(-2)(-3))+((-2)\cdot2-3\cdot3)\right|\\
&=|-12-5+0-13|=30,
\end{aligned}
\]
so the area is \(15\). The answer is \(\boxed{\textbf{(B) }15}\).
